\chapter{はじめに}
─内容─
参考文献（書籍）の引用\cite{book1}
\section{背景}
公立はこだて未来大学への通学手段は，公共交通機関において路線バスへの依存度が高い．特に，朝の登校ピーク時間帯には特定の便に乗客が集中し，車内の混雑が常態化している．この混雑は乗客の快適性を損なうだけでなく，乗降時間の長期化による運行ダイヤの遅延を引き起こし，学生の遅刻や一般利用者の利便性低下に悪影響を及ぼしている．また，バス事業者にとっても，定時運行およびサービス水準の維持が困難になるなど，運行管理上の問題となっている．

一方で，これらの問題に対する直接的な解決策である運行本数の増便は，近年のバス業界における運転手不足や運行コスト増加の観点から現実的には困難である．したがって，限られたリソースの中で運行効率を最大化するアプローチが求められている．

そこで本プロジェクトでは，既存の運行便のうち1本を主要停留所のみに停車する快速バスへと変更することで，乗客の乗車選択を分散させ，全体の混雑緩和および遅延抑制が可能であるという仮説を立て，研究テーマとして設定した．本報告では，プロジェクトの初期段階として，本テーマの決定に至った背景と解決すべき課題を整理するとともに，今後の検証に向けた到達目標について述べる．

\section{課題}
前項で述べたように，現在の公立はこだて未来大学へ向かう路線バスが抱える混雑および遅延の問題に対しては，単なるリソースの追加に頼るのではなく，既存の運行形態を見直し，乗客の利用便を分散させることで全体の最適化を図る具体的な方策を検討することが不可欠である．

しかし，実際の運行路線において，既存の便を快速化するなどのダイヤ変更を伴う実地検証を行うことは，利用者への影響や混乱のリスクが大きく，実環境での試行錯誤は困難である．

したがって，本研究においては，現在の需要動向や乗客の乗車選択行動を数理的にモデル化し，実社会に影響を与えることなく仮想空間上で新たな運行パターンの導入効果を定量的に評価・検証可能な手法を確立することが課題となっている．

\section{到達目標}
上記の課題を解決するため，本研究では数理モデリングを用いて既存の運行便のうち1便を快速化するシミュレーションを行い，その効果を評価・検証するアプローチを用いる．本研究の最終的な到達目標は，このアプローチによって，バスを利用する学生，地域住民，および運行事業者の3者すべてにメリットを生み出す最適な運行パターンの条件を定量的に明らかにすることである．

具体的には，大学への直行需要を快速便として分離することで，学生に対するピーク時の車内混雑を緩和し，通学の快適性と速達性を向上させる．また，学生の利用が快速便へ分散することにより，途中停留所から乗車する一般乗客の混雑による積み残しや車内空間の圧迫といった状況を改善し，地域公共交通としての確実な乗車機会と利便性を確保する．さらに，乗降時間の短縮によって運行ダイヤの遅延を抑制し，定時運行の達成を図る．これにより，限られたリソースのままで運送効率を最大化し，運行コストの最適化や信頼性向上を通じた運行事業者の利益確保に貢献する．